\section{Experimental methodology}
\label{sec:methodology}

\subsection{Research questions}
\begin{description}
\item[RQ1] By how much does HCA-DSE reduce the number of \emph{full design-point
  evaluations}, and how is that reduction distributed over the pruning layers?
\item[RQ2] Is the exact Pareto front preserved with respect to exhaustive enumeration?
\item[RQ3] At which evaluator cost does the pruning overhead pay for itself?
\item[RQ4] How does HCA-DSE compare with budget-matched stochastic exploration?
\item[RQ5] How do the reduction and the runtime scale with the size of the design space?
\item[RQ6] Do the architecture rankings produced by the analytical model transfer to a
  real containerised edge--cloud deployment?
\end{description}

\subsection{Design of the experiments}
Four design-space scales (S1--S4, Table~\ref{tab:variables}) are combined with the three
workload classes of Table~\ref{tab:workloads}. Exhaustive ground truth is computed for
S1--S3; S4 and the scalability study are explored by HCA-DSE only. Every configuration is
deterministic except the stochastic baselines, which are run with 20 independent seeds.

\paragraph{Baselines}
Exhaustive enumeration provides the ground-truth Pareto front. Random search and
NSGA-II~\cite{deb2002nsga2} (pymoo~\cite{blank2020pymoo}, population 40, SBX crossover,
polynomial mutation, integer repair) are run under matched evaluation budgets of
$100$, $250$, $500$ and $1000$ full evaluations, 20 seeds each. All methods use exactly
the same evaluator and the same feasibility rules, so differences are attributable to the
exploration strategy alone.

\paragraph{Metrics}
Pareto recall and precision are computed on sets of objective vectors against the
exhaustive front. Hypervolume~\cite{zitzler1999hv} is normalised by the hypervolume of the
exhaustive front (HV ratio); IGD$^{+}$~\cite{ishibuchi2015igdplus} uses the exhaustive
front as reference. We additionally report the number of full evaluations $N_E$, the
number of visited partial states $N$, and wall-clock time. Stochastic baselines are
compared with the Mann--Whitney $U$ test~\cite{mann1947whitney} and Cliff's
$\delta$~\cite{cliff1993delta}.

\paragraph{Ablation}
Four variants isolate the contribution of each layer: structural only; structural +
feasibility bounds; structural + dominance bounds; and the full method.

\subsection{Two-level evaluation and the cost model}
The analytical evaluator of Section~\ref{sec:model} (level A) is used for pruning,
ground truth and all exactness claims. A discrete-event simulator (level B) provides a
higher-fidelity, and therefore more expensive, evaluator: multi-server queues with a
finite number of servers per tier, explicit link contention, deterministic service times
and Poisson arrivals. Its measured cost is $19.4$\,ms per design point against
$\approx8\,\mu$s for the analytical evaluator, which allows the crossover analysis of
RQ3 to be anchored at two measured points rather than assumed.

\subsection{Testbed}
\label{sec:testbed}
The physical validation runs the architecture template as containers
(Fig.~\ref{fig:testbed}). Each edge node, gateway and cloud instance is a separate
container limited to one CPU, so that a node behaves as the single server assumed by the
model. Services are implemented in Go and share one binary: a request carries a plan
listing the remaining stages, and each node consumes CPU work equal to $w_i/C_t$ before
forwarding the remainder together with the real payload to the next hop. Links are shaped
per node with \texttt{tc netem}~\cite{hemminger2005netem} using the delay and rate of the
link class of the design under test; the gateway-to-cloud hop uses the WAN profile.
Replication is emulated faithfully: $r$ replicas execute stage $\tau_1$ in parallel and
carry the synchronisation payload, while the remainder of the chain is executed once.
A load generator produces Poisson arrivals at rate $\Lambda$, 900 requests per run of
which the first 150 are discarded as warm-up, and reports the latency percentiles.

Thirty-six architectures were selected across the three workloads --- Pareto-optimal,
dominated and near-boundary designs --- restricted to those that fit on a single machine
($\le4$ edge nodes, $\le2$ gateways, one cloud instance). Each was executed three times,
giving 108 measured runs.

\begin{figure}[t]
\centering
\begin{tikzpicture}[font=\scriptsize, node distance=5mm and 7mm,
  c/.style={draw, rounded corners=1.5pt, minimum height=6mm, minimum width=16mm,
            align=center, inner sep=2pt},
  arr/.style={-{Latex[length=1.8mm]}}]
\node[c, fill=black!5] (lg) {load generator\\[-2pt]Poisson $\Lambda$};
\node[c, right=of lg, fill=edgec!15, draw=edgec] (e1) {edge$_1$ \scriptsize 1 CPU};
\node[c, below=3mm of e1, fill=edgec!15, draw=edgec] (e2) {edge$_n$ \scriptsize 1 CPU};
\node[c, right=of e1, yshift=-4.5mm, fill=gwc!15, draw=gwc] (gw)
     {gateway \scriptsize 1 CPU};
\node[c, right=of gw, fill=cloudc!12, draw=cloudc] (cl) {cloud \scriptsize 1 CPU};
\draw[arr] (lg) -- (e1); \draw[arr] (lg) -- (e2);
\draw[arr] (e1) -- (gw); \draw[arr] (e2) -- (gw);
\draw[arr] (gw) -- node[above, font=\tiny] {WAN} (cl);
\node[font=\tiny, below=0.5mm of gw, text=black!60]
  {\texttt{tc netem}: per-node delay and rate};
\end{tikzpicture}
\caption{Containerised testbed. Every tier node is a separate one-CPU container; access
links carry the delay and bandwidth of the link class of the design, the cloud hop
carries the WAN profile.}
\label{fig:testbed}
\end{figure}

Because the analytical model is deliberately coarse, we report both the prediction error
(MAPE, median absolute error) and, more importantly for exploration, the \emph{ordering}
agreement: Spearman's $\rho$, Kendall's $\tau$, and the share of architecture pairs whose
measured order agrees with the model order among the pairs that the model separates by
more than $10\%$ and $25\%$.

\subsection{Reproducibility}
The implementation is in Python 3 with NumPy, SciPy and pymoo; the testbed services are
in Go and run under Docker Compose. All figures and all numeric tables of this paper are
generated from the raw result files by a single script, so no number is transcribed by
hand. Correctness is machine-checked at three levels: (i) bound admissibility is verified
\emph{exhaustively} on S1, i.e.\ for every prefix and every completion of that prefix
across the three workloads; (ii) for every pruned prefix, all of its completions are
enumerated and the claim of the corresponding proposition is verified directly, on S1 and
S2; and (iii) the equality $\mathcal{P}^{\mathrm{HCA}}=\PS$ is re-checked on S1--S3 for
every pruning variant as a regression test.
